% =========================================================
% Cauchy Sequences in Metric Spaces
% Source: volume-iii/analysis/metric-spaces/notes/cauchy-sequences.tex
%
% Dependencies:
%   notes-sequences.tex        — def:convergent-sequence, prop:limit-unique
%   notes-working-with-balls.tex — prop:reverse-triangle
% =========================================================

% ---------------------------------------------------------
% Toolkit
% ---------------------------------------------------------

\begin{tcolorbox}[colback=gray!6, colframe=gray!40, arc=2pt,
  left=6pt, right=6pt, top=4pt, bottom=4pt,
  title={\small\textbf{Cauchy Sequences --- Quick Reference}},
  fonttitle=\small\bfseries]
\small
\begin{tabular}{l l l l}
\toprule
\textbf{Label} & \textbf{Statement} & \textbf{Proof method} & \textbf{Detail} \\
\midrule
D\ref{def:cauchy-sequence}
  & $(x_n)$ Cauchy $\iff \forall\varepsilon>0,\;\exists N,\;m,n\ge N \Rightarrow d(x_m,x_n)<\varepsilon$
  & Definition
  & \hyperref[def:cauchy-sequence]{$\downarrow$} \\[4pt]
P\ref{prop:convergent-implies-cauchy}
  & Convergent $\Rightarrow$ Cauchy
  & Triangle inequality
  & \hyperref[prop:convergent-implies-cauchy]{$\downarrow$} \\[4pt]
P\ref{prop:cauchy-bounded}
  & Every Cauchy sequence is bounded
  & Finite initial segment
  & \hyperref[prop:cauchy-bounded]{$\downarrow$} \\[4pt]
P\ref{prop:cauchy-subsequence}
  & Cauchy $+$ convergent subsequence $\Rightarrow$ convergent
  & Triangle inequality
  & \hyperref[prop:cauchy-subsequence]{$\downarrow$} \\[4pt]
P\ref{prop:cauchy-unique-limit}
  & Cauchy sequence has at most one limit
  & Uniqueness of limits
  & \hyperref[prop:cauchy-unique-limit]{$\downarrow$} \\[4pt]
\bottomrule
\end{tabular}
\end{tcolorbox}

% ---------------------------------------------------------
\subsection{The Cauchy Condition}
\label{sec:cauchy-condition}
% ---------------------------------------------------------

Recall that a sequence $(x_n)$ in a metric space $(X,d)$ converges to a
point $x \in X$ if for every $\varepsilon > 0$ there exists $N$ such that
$n \ge N$ implies $d(x_n, x) < \varepsilon$.
This is an \emph{external} condition: it tests the sequence against a
specific candidate limit $x$.

The Cauchy condition removes all reference to a limit.
It asks only whether the terms of the sequence are eventually close to
\emph{each other}.

\begin{tcolorbox}[colback=propbox, colframe=propborder, arc=2pt,
  left=6pt, right=6pt, top=4pt, bottom=4pt,
  title={\small\textbf{Definition (Cauchy Sequence)}},
  fonttitle=\small\bfseries]
\begin{definition}[Cauchy Sequence]
\label{def:cauchy-sequence}
Let $(X,d)$ be a metric space.
A sequence $(x_n)_{n\in\mathbb{N}}$ in $X$ is a \emph{Cauchy sequence} if
\[
  \forall \varepsilon > 0,\;
  \exists N \in \mathbb{N},\;
  \forall m,n \ge N,\quad
  d(x_m, x_n) < \varepsilon.
\]
\end{definition}
\end{tcolorbox}

\begin{remark*}[Fully quantified form]
\[
(x_n)\ \text{Cauchy}
\;\iff\;
\forall \varepsilon \in \mathbb{R}^{+},\;
\exists N \in \mathbb{N},\;
\forall m \in \mathbb{N},\;
\forall n \in \mathbb{N},\;
\bigl(m \ge N \;\land\; n \ge N\bigr)
\;\Rightarrow\;
d(x_m, x_n) < \varepsilon.
\]
\end{remark*}

\begin{remark*}[Intuition]
The terms of a Cauchy sequence ``bunch together'' as $n \to \infty$.
For any prescribed tolerance $\varepsilon$, all sufficiently late
pairs of terms lie within $\varepsilon$ of each other.
No external reference point is involved: the condition is entirely
about the mutual spacing of the sequence's own terms.
\end{remark*}

\begin{remark*}[Comparison with convergence]
\begin{center}
\begin{tabular}{l l}
\toprule
\textbf{Concept} & \textbf{Condition} \\
\midrule
$(x_n) \to x$ (convergent)
  & $\forall\varepsilon>0,\;\exists N,\;\forall n\ge N:\;
    d(x_n,x) < \varepsilon$ \\[4pt]
$(x_n)$ Cauchy
  & $\forall\varepsilon>0,\;\exists N,\;\forall m,n\ge N:\;
    d(x_m,x_n) < \varepsilon$ \\
\bottomrule
\end{tabular}
\end{center}
\smallskip\noindent
Convergence involves one index ranging over a tail; the Cauchy condition
involves \emph{two} independent indices both ranging over the same tail.
\end{remark*}

\begin{remark*}[Common error: consecutive spacing is not enough]
The condition $d(x_{n+1}, x_n) \to 0$ --- that \emph{consecutive} terms
approach each other --- is strictly weaker than the Cauchy condition.

The harmonic series partial sums $s_n = \sum_{k=1}^n \frac{1}{k}$ in
$\mathbb{R}$ satisfy $d(s_{n+1}, s_n) = \frac{1}{n+1} \to 0$, yet the
sequence diverges, so it is not Cauchy.

The Cauchy condition controls the distance between \emph{all} pairs
$x_m, x_n$ with $m,n \ge N$, not just adjacent ones.
\end{remark*}

% ---------------------------------------------------------
\subsection{Basic Properties}
\label{sec:cauchy-basic-properties}
% ---------------------------------------------------------

\begin{proposition}[\texorpdfstring{\hyperref[prf:convergent-implies-cauchy]{Convergent Implies Cauchy}}{Convergent Implies Cauchy}]
\label{prop:convergent-implies-cauchy}
Let $(X,d)$ be a metric space.
Every convergent sequence in $X$ is a Cauchy sequence.
\end{proposition}

\begin{proof}
Let $(x_n) \to x$ in $X$.
Given $\varepsilon > 0$, choose $N \in \mathbb{N}$ so that for all
$n \ge N$, $d(x_n, x) < \varepsilon/2$.
For any $m, n \ge N$, the triangle inequality gives
\[
  d(x_m, x_n)
  \;\le\;
  d(x_m, x) + d(x, x_n)
  \;<\;
  \tfrac{\varepsilon}{2} + \tfrac{\varepsilon}{2}
  \;=\;
  \varepsilon. \qedhere
\]
\end{proof}

\begin{remark*}[The converse]
The converse --- every Cauchy sequence converges --- fails in general
metric spaces and is precisely the property of \emph{completeness}.
See \hyperref[sec:complete-metric-spaces]{\S\ref*{sec:complete-metric-spaces}}.
In $\mathbb{Q}$: the sequence of rational truncations
$1,\; 1.4,\; 1.41,\; 1.414, \ldots$
of $\sqrt{2}$ is Cauchy but has no limit in $\mathbb{Q}$.
\end{remark*}

\begin{proposition}[\texorpdfstring{\hyperref[prf:cauchy-bounded]{Cauchy Sequences are Bounded}}{Cauchy Sequences are Bounded}]
\label{prop:cauchy-bounded}
Let $(X,d)$ be a metric space.
Every Cauchy sequence in $X$ is bounded.
\end{proposition}

\begin{proof}
Let $(x_n)$ be Cauchy.
Apply the definition with $\varepsilon = 1$: there exists $N \in \mathbb{N}$
such that for all $m,n \ge N$, $d(x_m,x_n) < 1$.
In particular, $d(x_n, x_N) < 1$ for all $n \ge N$.

The finite set of initial terms contributes at most
\[
  R
  \;:=\;
  \max\bigl\{
    d(x_1, x_N),\;
    d(x_2, x_N),\;
    \ldots,\;
    d(x_{N-1}, x_N),\;
    1
  \bigr\}.
\]
Then $d(x_n, x_N) \le R$ for all $n \in \mathbb{N}$, so every
term lies in $\overline{B}_R(x_N)$.
\end{proof}

\begin{remark*}[Fully quantified form]
\[
\forall (X,d)\ \text{metric space},\;
\forall (x_n)\ \text{Cauchy in } X:\quad
\exists x_0 \in X,\;
\exists R \in \mathbb{R}^{+},\;
\forall n \in \mathbb{N},\;
d(x_n, x_0) \le R.
\]
\end{remark*}

\begin{remark*}[Consequence]
Boundedness is a necessary condition for a sequence to be Cauchy.
Any unbounded sequence is therefore not Cauchy, which gives a quick
way to rule out the Cauchy property in examples.
\end{remark*}

\begin{proposition}[\texorpdfstring{\hyperref[prf:cauchy-subsequence]{Cauchy with Convergent Subsequence Converges}}{Cauchy with Convergent Subsequence Converges}]
\label{prop:cauchy-subsequence}
Let $(X,d)$ be a metric space.
If $(x_n)$ is a Cauchy sequence and $(x_{n_k})$ is a subsequence that
converges to some $x \in X$, then $(x_n) \to x$.
\end{proposition}

\begin{proof}
Given $\varepsilon > 0$:
\begin{enumerate}[label=(\roman*), leftmargin=2em]
  \item Since $(x_n)$ is Cauchy, choose $N_1$ so that
        $m,n \ge N_1$ implies $d(x_m,x_n) < \varepsilon/2$.
  \item Since $x_{n_k} \to x$, choose $K$ so that $k \ge K$ implies
        $d(x_{n_k}, x) < \varepsilon/2$.
\end{enumerate}
Pick any index $k \ge K$ with $n_k \ge N_1$.
For any $n \ge N_1$,
\[
  d(x_n, x)
  \;\le\;
  d(x_n, x_{n_k}) + d(x_{n_k}, x)
  \;<\;
  \tfrac{\varepsilon}{2} + \tfrac{\varepsilon}{2}
  \;=\;
  \varepsilon. \qedhere
\]
\end{proof}

\begin{remark*}[Role of this proposition]
This is the key technical bridge in most proofs that a specific space
is complete.
The strategy is:
\begin{enumerate}[label=(\arabic*), leftmargin=2em]
  \item Take an arbitrary Cauchy sequence $(x_n)$.
  \item Use boundedness (Proposition~\ref{prop:cauchy-bounded}) together
        with a selection theorem (e.g.\ Bolzano--Weierstrass in
        $\mathbb{R}^n$) to extract a convergent subsequence.
  \item Apply this proposition to conclude the full sequence converges.
\end{enumerate}
\end{remark*}

\begin{proposition}[\texorpdfstring{\hyperref[prf:cauchy-unique-limit]{Cauchy Sequences Have at Most One Limit}}{Cauchy Sequences Have at Most One Limit}]
\label{prop:cauchy-unique-limit}
Let $(X,d)$ be a metric space.
If $(x_n)$ is a Cauchy sequence that converges to both $x$ and $y$ in $X$,
then $x = y$.
\end{proposition}

\begin{proof}
This follows immediately from the uniqueness of limits in metric spaces.
For completeness: given $\varepsilon > 0$, choose $N$ so that
$n \ge N$ implies $d(x_n,x) < \varepsilon/2$ and $d(x_n,y) < \varepsilon/2$.
Then
\[
  d(x,y) \;\le\; d(x, x_N) + d(x_N, y) \;<\; \varepsilon.
\]
Since $\varepsilon$ was arbitrary, $d(x,y) = 0$, hence $x = y$.
\end{proof}

% ---------------------------------------------------------
\subsection{Examples and Non-Examples}
\label{sec:cauchy-examples}
% ---------------------------------------------------------

\begin{example}[Cauchy in $\mathbb{R}$]
The sequence $x_n = \frac{1}{n}$ is Cauchy in $\mathbb{R}$.
Given $\varepsilon > 0$, choose $N > 2/\varepsilon$.
For $m, n \ge N$:
\[
  d(x_m, x_n)
  = \Bigl|\frac{1}{m} - \frac{1}{n}\Bigr|
  \le \frac{1}{m} + \frac{1}{n}
  \le \frac{2}{N}
  < \varepsilon.
\]
This sequence converges to $0 \in \mathbb{R}$, confirming
Proposition~\ref{prop:convergent-implies-cauchy}.
\end{example}

\begin{example}[Cauchy in $\mathbb{Q}$ without a rational limit]
Let $x_n$ be the decimal truncation of $\sqrt{2}$ to $n$ decimal places:
\[
  x_1 = 1.4,\quad
  x_2 = 1.41,\quad
  x_3 = 1.414,\quad
  x_4 = 1.4142,\quad \ldots
\]
For $m \ge n$, $|x_m - x_n| \le 10^{-n}$, so given $\varepsilon > 0$
choose $N > -\log_{10}\varepsilon$.
Then $m,n \ge N$ gives $|x_m - x_n| \le 10^{-N} < \varepsilon$.

The sequence is Cauchy in $\mathbb{Q}$.
In $\mathbb{R}$ it converges to $\sqrt{2}$, but $\sqrt{2} \notin \mathbb{Q}$,
so this Cauchy sequence has no limit in $\mathbb{Q}$.
This is the canonical example of an incomplete metric space.
\end{example}

\begin{example}[Not Cauchy: consecutive spacing decays but sequence diverges]
In $\mathbb{R}$, let $x_n = \sum_{k=1}^{n} \frac{1}{k}$ (partial sums of the
harmonic series).
Then $d(x_{n+1}, x_n) = \frac{1}{n+1} \to 0$, yet for $m = 2n$:
\[
  d(x_{2n}, x_n)
  = \sum_{k=n+1}^{2n} \frac{1}{k}
  \ge \sum_{k=n+1}^{2n} \frac{1}{2n}
  = \frac{1}{2}.
\]
No matter how large $N$ is, taking $n = N$ and $m = 2N$ gives a pair
with $d(x_m, x_n) \ge 1/2$.
So $(x_n)$ is not Cauchy, despite consecutive terms becoming arbitrarily
close.
\end{example}

\begin{example}[Discrete metric: Cauchy sequences are eventually constant]
In any metric space with the discrete metric
$d(x,y) = \mathbf{1}[x \neq y]$,
a sequence $(x_n)$ is Cauchy if and only if it is eventually constant.

\emph{Proof.} Apply the definition with $\varepsilon = 1/2$: there
exists $N$ such that for all $m,n \ge N$, $d(x_m, x_n) < 1/2$.
Since the discrete metric takes values in $\{0,1\}$, this forces
$d(x_m,x_n) = 0$, i.e.\ $x_m = x_n$ for all $m,n \ge N$.
\end{example}

% ---------------------------------------------------------
\subsection{The Cauchy Condition and Subsequences}
\label{sec:cauchy-subsequences}
% ---------------------------------------------------------

\begin{proposition}[Subsequences of Cauchy Sequences are Cauchy]
\label{prop:cauchy-subsequence-cauchy}
Let $(X,d)$ be a metric space.
Every subsequence of a Cauchy sequence is Cauchy.
\end{proposition}

\begin{proof}
Let $(x_n)$ be Cauchy and let $(x_{n_k})$ be a subsequence.
Since $n_k \ge k$ for all $k$, given $\varepsilon > 0$ choose $N$ for the
Cauchy condition of $(x_n)$.
For $j, k \ge N$, we have $n_j, n_k \ge N$, so
$d(x_{n_j}, x_{n_k}) < \varepsilon$.
\end{proof}

\begin{remark*}[Consequence for the completion construction]
This proposition is used in the proof of the Completion Theorem: given a
Cauchy sequence of equivalence classes in the completion, one selects
convenient representatives by passing to subsequences.
Since any subsequence of a Cauchy sequence is again Cauchy, the
representatives remain in the right category.
See \hyperref[sec:completion]{\S\ref*{sec:completion}}.
\end{remark*}

% ---------------------------------------------------------
\subsection{Conceptual Summary}
\label{sec:cauchy-summary}
% ---------------------------------------------------------

\begin{center}
\begin{tabular}{l l}
\toprule
\textbf{Property} & \textbf{Status} \\
\midrule
Convergent $\Rightarrow$ Cauchy & Always (Prop.~\ref{prop:convergent-implies-cauchy}) \\
Cauchy $\Rightarrow$ Convergent & Only in complete spaces \\
Cauchy $\Rightarrow$ Bounded    & Always (Prop.~\ref{prop:cauchy-bounded}) \\
Cauchy $+$ convergent subsequence $\Rightarrow$ convergent
  & Always (Prop.~\ref{prop:cauchy-subsequence}) \\
Cauchy has at most one limit    & Always (Prop.~\ref{prop:cauchy-unique-limit}) \\
Subsequence of Cauchy is Cauchy & Always (Prop.~\ref{prop:cauchy-subsequence-cauchy}) \\
Consecutive spacing $\to 0$ $\Rightarrow$ Cauchy & \emph{False in general} \\
\bottomrule
\end{tabular}
\end{center}

\begin{center}
\[
  \boxed{
    \text{Cauchy} = \text{``terms are mutually close in the tail''}
  }
\]
\[
  \boxed{
    \text{Whether Cauchy} \Rightarrow \text{convergent is the question of completeness.}
  }
\]
\end{center}

\begin{remark*}[Forward reference]
Whether every Cauchy sequence in a given metric space converges --- the
property of \emph{completeness} --- is the subject of
\hyperref[sec:complete-metric-spaces]{\S\ref*{sec:complete-metric-spaces}}.
\end{remark*}